\section{Conclusion}\label{sec-conc}

This paper proposes a large series prior for Bayesian nonparametric regression: a deterministic number of basis functions growing as $n^b$, combined with a shrinkage prior on the coefficients. Theorem~\ref{thm:main} gives posterior contraction rates for product shrinkage priors that satisfy prior-mass conditions adapted from \citet{WeiGhosal2020}, with general dictionaries, and shows that the posterior concentrates on coefficient vectors with at most a multiple of $s_n$ coordinates above the negligibility threshold.

Under the sparse approximation profile and the prior conditions of Corollary~\ref{cor:anisotropic}, the posterior contracts at the rate $(n/\log n)^{-\alpha/(2\alpha+1)}$ for $f_0\in C^\alpha([0,1])$, the minimax rate up to a logarithmic factor, and concentrates on coefficient vectors with at most a fixed multiple of $(n/\log n)^{1/(2\alpha+1)}$ coordinates exceeding $\delta_n$ in absolute value. In $d$ dimensions the rate is $(n/\log n)^{-\alpha/(2\alpha+d)}$ and the support bound is a fixed multiple of $(n/\log n)^{d/(2\alpha+d)}$; under anisotropic smoothness, $\alpha$ is replaced by the harmonic mean $\alpha^*$ of the directional smoothness parameters. These conclusions apply to the classes of Lemmas~\ref{lem:bspline-profile} and~\ref{lem:fourier-profile}: periodic functions for the Fourier dictionary, and smoothness at most the spline order, with integer directional smoothness in the anisotropic case, for the multiresolution B-spline dictionary. The Horseshoe with a deterministic global scale satisfies the prior conditions, and its scale does not depend on the smoothness, so one prior attains these rates pointwise over the admissible range. More broadly, the rate is governed by the sparsity with which the dictionary represents $f_0$ (Remark~\ref{rem:polynomial}), so a basis system matters mainly through the sparsity it induces.

The numerical results support this picture. On a common Fourier dictionary, the Horseshoe has a lower error than cross-validated ridge regression in every comparison, and its error stays essentially flat as the dictionary grows, while the ridge error roughly doubles. In the scaling experiment, the empirical error decay follows the sparsity of the fitted coefficients, and the basis system, rather than the smoothness of the target, separates the rates. In the three real-data applications, one rule for the dictionary size and the global scale yields smooth fits of very different complexity. A natural next step is a sampler that remains stable at the very small global scales the theory calls for, which would bring the simulations closer to the regime of Theorem~\ref{thm:main}.
